\subsection{A multiplicity-space obstruction to the isotypic projector extension}

\begin{proposition}[P-F61: whole irreducible channels do not determine the nonstandard sign projector]
Let $n\ge3$, $H=GL_n(\mathbb C)$, and
\[
 X=\mathfrak{sl}_n\otimes\det
       \simeq S^{(2,1^{n-2})}\mathbb C^n.
                                                               \tag{P-F61a}
\]
Both $\operatorname{Sym}^2X$ and $\bigwedge^2X$ contain the
same irreducible $H$-module
\[
 Y=\mathfrak{sl}_n\otimes\det^2
       \simeq S^{(3,2^{n-2},1)}\mathbb C^n.
                                                               \tag{P-F61b}
\]
Consequently the alternating projector
$Q^-=(1-\tau)/2$ on $X\otimes X$, where $\tau$ exchanges the
factors, is not central in $\operatorname{End}_H(X\otimes X)$.
It cannot be any linear combination of the projections onto
whole irreducible isotypic components.

In particular the following proposed extension of the
$\mathfrak{sl}_2$ channel construction is impossible: represent
the nonstandard negative-braid projector by a combination of
whole isotypic projectors whose coefficients and projectors
are regular at $q=1$ and whose limits are the corresponding
classical isotypic projectors.  This excludes that specified
isotypic construction, not categorifications using finer
multiplicity-space projectors.
\end{proposition}

\begin{proof}
Conjugation makes $\mathfrak{sl}_n$ the irreducible adjoint
$SL_n$-module.  Its $GL_n$ highest weight is
$(1,0^{n-2},-1)$, represented by the matrix unit $E_{1n}$.
Tensoring by $\det$ and $\det^2$ adds one and two to every
entry, respectively, proving the polynomial-module
identifications in (P-F61a)--(P-F61b).

For traceless matrices define
\[
 B(A,B)=AB-BA,\qquad
 J(A,B)=AB+BA-\frac2n\operatorname{tr}(AB)I_n.
                                                               \tag{P-F61c}
\]
The first bilinear map is alternating and the second is
symmetric.  Both are conjugation-equivariant and take values
in $\mathfrak{sl}_n$.  The bracket is nonzero because
$B(E_{12},E_{23})=E_{13}$.
For $H_0=\operatorname{diag}(1,-1,0,\ldots,0)$,
\[
 J(H_0,H_0)=2\operatorname{diag}(1,1,0,\ldots,0)
                            -\frac4nI_n\ne0
          \qquad(n\ge3).
                                                               \tag{P-F61d}
\]
Each image is a nonzero submodule of the irreducible adjoint
module, hence is the entire adjoint module.  After tensoring
with $\det^2$ this gives surjections
\[
 \bigwedge^2X\longrightarrow Y,
       \qquad \operatorname{Sym}^2X\longrightarrow Y.
                                                               \tag{P-F61e}
\]
Finite-dimensional rational representations of $GL_n$ over
$\mathbb C$ are completely reducible.  Splitting these
surjections gives submodules $Y_-\subset\bigwedge^2X$ and
$Y_+\subset\operatorname{Sym}^2X$, each isomorphic to $Y$.
This uses only complete reducibility; no choice of the
splittings is claimed to be canonical.

Choose an equivariant isomorphism $T:Y_+\to Y_-$ and extend
it by zero on an equivariant complement of $Y_+$ in
$X\otimes X$.  On $Y_+$ one has
\[
           (Q^-T-TQ^-)\big|_{Y_+}=T\ne0,
                                                               \tag{P-F61f}
\]
since $Q^-$ is zero on the symmetric square and the identity
on the exterior square.  Thus $Q^-$ is not central.  Equivalently,
on the $Y$-isotypic component every whole-isotypic projector
is scalar, whereas $Q^-$ has both eigenvalues zero and one.
This also proves directly that no linear combination of
such projectors equals $Q^-$.

For the stated quantum extension, take the limit at $q=1$
of the proposed identity.  By its regularity and specialization
hypotheses the right side is a linear combination of the
classical whole-isotypic projectors.  The left side is $Q^-$:
the braid specializes to the flip, and its positive and negative
spectral groups specialize to the $+1$ and $-1$ flip eigenspaces.
The latter groups remain separated near $q=1$, so their
spectral projectors have these limits.  This contradicts
(P-F61f), proving precisely the asserted failure.
\end{proof}

\paragraph{Scope of the categorical sources and the remaining construction.}
The source nonstandard projector is the sign projection of the
actual braid, GCT VII, Section~2, equations~(4)--(8), printed
pp.~6--7; its two signs, not merely a total highest-weight
label, define the algebra.  Cautis, \emph{Clasp technology to
knot homology via the affine Grassmannian}, arXiv:1207.2074v3,
Theorems~2.2--2.3, printed pp.~6--7, constructs infinite-twist
clasps projecting a tensor product of fundamental modules onto
its Cartan highest-weight summand.  Theorem~2.6, p.~7, realizes
the required action using derived coherent sheaves on affine
Grassmannian convolution varieties:
\url{https://arxiv.org/pdf/1207.2074}.
These are actual geometric categorical inputs, but their cited
statements do not identify the sign split inside (P-F61b).

Elias--Hogancamp, \emph{Categorical diagonalization of full
twists}, arXiv:1801.00191v2, Theorem~1.7, printed pp.~7--8,
constructs both central and primitive tableau projectors in
the Hecke setting.  Its primitive projectors are finer than
whole-isotypic projectors and are not excluded by P-F61.
Their theorem uses a tower of full twists; the paragraph following
Remark~1.6, p.~7, explains why a single full twist does not
separate all the required labels.  Remark~1.8, p.~8, distinguishes the precise
centrality proved there from stronger coherent centrality:
\url{https://arxiv.org/pdf/1801.00191}.
To use such finer projectors here one must identify their
actual corner and multiplicity-space action with the
nonstandard braid-sign projector, and then construct the
required compatible mixed relations.  P-F61 refutes the
whole-isotypic shortcut, leaving that more precise comparison
as the next higher-rank step.  It does not refute any of the
cited clasp or tableau-projector theorems, and makes no claim
of historical priority for the classical adjoint calculation.
